\section{Methodology: Causal MARL for Equilibrium Robustness}

In this section, we transition from the conceptual challenge of non-stationarity to a formal mathematical framework that treats decentralized agents not as mere reward-optimizers, but as causal reasoners. We first define the Decentralized Causal Markov Decision Process (D-CMDP) as our foundational structure, then derive the conditions for invariant causal discovery, and finally present the \textsc{Causal-Equilibrium} algorithm.

\subsection{Formalizing Causal Awareness in Decentralized Systems}

Traditional MARL typically operates within the framework of a Decentralized Partially Observable Markov Decision Process (Dec-POMDP) \cite{bernstein2002complexity}. However, these models often ignore the structural dependencies between agent actions and environmental transitions, treating other agents' policies as exogenous noise. We extend this by introducing the \textit{Decentralized Causal Sequential Decision-Making} framework.

Let $\mathcal{M} = \langle \mathcal{S}, \{\mathcal{A}_i\}_{i=1}^N, \mathcal{T}, \{R_i\}_{i=1}^N, \gamma, \mathcal{G} \rangle$ be a D-CMDP, where $\mathcal{G}$ is a Structural Causal Model (SCM) \cite{pearl2009causality} representing the instantaneous and temporal causal dependencies. For each time step $t$, the state transition $s_{t+1}$ and rewards $r_{i,t}$ are governed by structural equations:
\begin{equation}
s_{t+1} = f_{\mathcal{S}}(s_t, a_{1,t}, \dots, a_{N,t}, \epsilon_{s,t}), \quad r_{i,t} = f_{R_i}(s_t, a_{i,t}, a_{-i,t}, \epsilon_{r,t})
\end{equation}
where $a_{-i}$ denotes the actions of all agents except $i$, and $\epsilon$ represents independent exogenous noise.

\textbf{Definition 1 (Causal Agent).} An agent $i$ is \textit{causally aware} if it maintains a partial estimate $\hat{\mathcal{G}}_i$ of the SCM and utilizes the $do$-operator to evaluate counterfactual policies. Specifically, the agent optimizes a policy $\pi_i$ by considering the interventional distribution $P(s_{t+1} | s_t, do(a_{i,t}), \pi_{-i})$, rather than the mere conditional distribution $P(s_{t+1} | s_t, a_{i,t}, \pi_{-i})$ which is prone to spurious correlations induced by the latent strategy shifts of other agents.

This distinction is critical: in high-frequency environments like real-time bidding or algorithmic trading, ``spurious equilibria'' often emerge because agents mistake the shifting distribution of competitor strategies for environmental dynamics. Causal awareness allows the agent to disentangle the \textit{invariant} laws of the system from the \textit{transient} behaviors of competitors.

\subsection{Causal Discovery of Invariant Relationships}

To achieve robustness, agents must identify causal links that are invariant to the non-stationary policies of other agents \cite{peters2016causal}. We define the \textit{Invariant Causal Set} $\mathcal{C}^* \subseteq \mathcal{S} \cup \mathcal{A}$ as the set of variables that have a stable causal effect on agent $i$'s long-term utility across a set of environments $\mathcal{E}$ (representing different competitor behaviors).

\textbf{The Invariance Principle in MARL:} A causal relationship $a_{i,t} \to s_{t+1}$ is invariant if for any two joint policies $\pi_{-i}, \pi'_{-i}$ of other agents, the structural mechanism remains unchanged:
\begin{equation}
P(s_{t+1} | \text{pa}(s_{t+1}), do(a_{i,t}))_{\pi_{-i}} = P(s_{t+1} | \text{pa}(s_{t+1}), do(a_{i,t}))_{\pi'_{-i}}
\end{equation}
where $\text{pa}(s_{t+1})$ are the parents of $s_{t+1}$ in $\mathcal{G}$. We utilize a variation of the PC algorithm \cite{spirtes2000causation} adapted for time-series, which we call \textit{Multi-Agent Interventional PC (MAI-PC)}. Unlike standard discovery, MAI-PC uses the agents' own exploratory actions as natural interventions to test for conditional independence. 

By prioritizing the discovery of these invariant links, our framework avoids the ``Latency Illusion''---the tendency of agents to overfit to short-term, coincidental patterns in the competitive landscape. This leads directly into the derivation of stable equilibrium properties.

\subsection{Equilibrium Properties under Causal Awareness}

We now introduce the concept of the \textbf{Causal Nash Equilibrium (CNE)}. In a CNE, each agent $i$ plays a policy $\pi^*_i$ that is a best response to the interventional distributions generated by $\pi^*_{-i}$, specifically robust to distributional shifts in the non-invariant parts of the system.

\textbf{Theorem 1 (Robustness of CNE).} \textit{Let $\mathcal{G}$ be the true causal graph of the system. If agents reach an equilibrium where each agent $i$ minimizes the Causal Risk $J_i(\pi_i) = \mathbb{E}_{\pi_i, \pi_{-i}, e \in \mathcal{E}} [ \sum \gamma^t R_{i,t} ]$ subject to the constraint of invariance across $\mathcal{E}$, then the resulting equilibrium is stable against any $\delta$-perturbation in the strategies of a sub-coalition of agents $j \notin i$, provided $\delta$ does not alter the underlying structural equations.}

This theorem provides the ``Why'' behind our approach: causal awareness serves as a regularization mechanism that prevents the system from drifting into fragile states. In non-ergodic environments, where historical averages do not predict future states, the CNE relies on the \textit{physics} of the interaction rather than the \textit{statistics} of the past \cite{eichler2012causal, peters2017elements}.

\subsection{Learning Invariant Causal Representations}

Having established the theory, we present the algorithmic implementation: \textsc{Causal-Equilibrium-Net} (CE-Net). The architecture consists of three components: (1) a \textit{Causal Encoder} that maps high-dimensional observations to a latent SCM, (2) an \textit{Invariance Discriminator}, and (3) a \textit{Causal Policy Head}.

\begin{algorithm}[h]
\caption{\textsc{Causal-Equilibrium} (CE) Training Loop}
\begin{algorithmic}[1]
\STATE \textbf{Initialize:} Causal Graph $\hat{\mathcal{G}}$, Policy $\pi_i$, Encoder $\phi$, Buffer $\mathcal{D}$
\FOR{episode $e = 1$ \TO $M$}
    \STATE Sample trajectory $\tau = \{s_t, a_t, r_t\}$ using current policies
    \STATE \textit{// Causal Discovery Step}
    \STATE Update $\hat{\mathcal{G}}$ using MAI-PC on $\mathcal{D}$ with interventional tests \cite{pearl2009causality}
    \STATE \textit{// Representation Learning}
    \STATE $\mathcal{L}_{inv} = \sum_{j \in \text{Environments}} \| P(\phi(s_{t+1}) | \phi(s_t), a_{i,t})_{\pi_{-i, j}} - \bar{P} \|^2$ \cite{arjovsky2019invariant}
    \STATE \textit{// Policy Optimization}
    \STATE Calculate Causal Gradient: $\nabla_{\theta} J(\theta) = \mathbb{E} [ \nabla_{\theta} \log \pi(a|s) \cdot \hat{Q}^{\text{causal}}(s, a) ]$
    \STATE Update $\theta$ using Adam optimizer \cite{kingma2014adam}
\ENDFOR
\end{algorithmic}
\end{algorithm}

The Causal Encoder minimizes a multi-objective loss function:
\begin{equation}
\mathcal{L} = \mathcal{L}_{reconstruct} + \lambda_1 \mathcal{L}_{sparsity}(\hat{\mathcal{G}}) + \lambda_2 \mathcal{L}_{invariance}(\Phi)
\end{equation}
where $\mathcal{L}_{invariance}$ is inspired by Invariant Risk Minimization (IRM) \cite{arjovsky2019invariant, zhang2024causal}, ensuring that the features used for decision-making have a stable relationship with the reward across different competitive ``environments'' or strategy regimes.

\subsection{Causal Policy Optimization and Complexity}

The final component is the \textit{Causal Policy Gradient}. Standard policy gradients $\nabla \mathbb{E}[R]$ are often high-variance in multi-agent settings due to the ``moving target'' problem. We reformulate the gradient using the \textit{Causal Jacobian}:
\begin{equation}
g_{causal} = \mathbb{E}_{s \sim \rho^\pi} \left[ \frac{\partial \pi(a|s)}{\partial \theta} \left( \sum_{s'} \frac{\partial P(s' | s, do(a))}{\partial a} \mathcal{V}(s') \right) \right]
\end{equation}
where $\frac{\partial P(s' | s, do(a))}{\partial a}$ represents the sensitivity of the environment to the agent's actions, derived from the learned SCM \cite{wang2024causal_rl}. This gradient provides a more direct path to stable regions of the strategy space.

\textbf{Computational Complexity:} While causal discovery is NP-hard in the general case, our use of MAI-PC on a restricted latent space $\mathcal{Z} \ll \mathcal{S}$ reduces the complexity to $O(N \cdot |\mathcal{Z}|^k)$, where $k$ is the maximum in-degree of the causal graph. This allows our agents to perform discovery in real-time within high-frequency loops, effectively overcoming the ``Latency Illusion'' through architectural efficiency rather than brute-force computation.

By integrating these components, \textsc{Causal-Equilibrium} moves the system toward a state where stability is a structural property of the agents' understanding, rather than a coincidental outcome of optimization. This paves the way for our empirical evaluation in complex, non-ergodic environments.